\section{Results}

\subsection{Strict leave-one-city-out dynamics}

\Cref{fig:strict-loo,tab:strict-loo} show the primary dynamics result.
Unweighted \ours beats the dense \Htwo-style residual baseline on all four held-out cities, with mean total-MSE ratio $0.759$.
Similarity-weighted \ours also wins on all four targets and improves the mean ratio to $0.578$.
The largest gains occur for Austin, Halifax, and MBTA.
Singapore is the hardest target: unweighted \ours is only slightly better than \Htwo, and similarity weighting is slightly worse than unweighted \ours while still remaining below the \Htwo baseline.

\begin{figure}[t]
\centering
\includegraphics[width=0.78\textwidth]{strict_leave_one_out_total_mse_ratio.png}
\caption{Strict leave-one-city-out dynamics validation. Bars show total-MSE ratio relative to \Htwo. The dashed line is parity with \Htwo.}
\label{fig:strict-loo}
\end{figure}

\begin{table}[t]
\centering
\scriptsize
\resizebox{\textwidth}{!}{\begin{tabular}{lllllllll}
\toprule
target\_city & source\_cities & target\_transitions & h2oplus\_total\_mse & cfcmt\_total\_mse & weighted\_cfcmt\_total\_mse & cfcmt\_vs\_h2oplus\_ratio & weighted\_cfcmt\_vs\_h2oplus\_ratio & weighted\_cfcmt\_vs\_unweighted\_ratio \\
\midrule
Singapore & austin\_capmetro\_all, halifax\_transit\_all, mbta\_all & 1243488 & 2254.8548 & 2246.6693 & 2251.7081 & 0.9964 & 0.9986 & 1.0022 \\
Austin / CapMetro & singapore\_lta\_all, halifax\_transit\_all, mbta\_all & 310272 & 91.9724 & 67.8254 & 36.6531 & 0.7375 & 0.3985 & 0.5404 \\
Halifax Transit & singapore\_lta\_all, austin\_capmetro\_all, mbta\_all & 266448 & 97.1175 & 69.5992 & 54.7536 & 0.7166 & 0.5638 & 0.7867 \\
MBTA & singapore\_lta\_all, austin\_capmetro\_all, halifax\_transit\_all & 969360 & 318.9925 & 186.4177 & 112.2698 & 0.5844 & 0.352 & 0.6022 \\
\bottomrule
\end{tabular}
}
\caption{Strict leave-one-city-out dynamics results. Ratios below 1 indicate lower total MSE than \Htwo.}
\label{tab:strict-loo}
\end{table}

\subsection{Target-information ladder}

\Cref{tab:target-information-ladder} separates the regimes that are easy to confuse in cross-city transit transfer.
The single-source zero-shot diagnostic is deliberately brittle: across the 12 ordered source-target pairs, \ours wins against \Htwo on 6 pairs with mean ratio $0.927$.
This shows that ``train on one city, deploy in another'' is possible in some directions but not a stable main claim.
Adding source diversity changes the result: the multi-source target-label-free setting wins on all four held-out cities with mean weighted ratio $0.578$.
With small target-route historical labels, the same framework becomes target offline adaptation rather than zero-shot transfer; the $1\%$--$25\%$ route-budget points remain below \Htwo-style dense adaptation on three of four strict targets.
The $100\%$ target-route point is kept only as an upper-bound reference.

\begin{center}
\centering
\scriptsize
\resizebox{\textwidth}{!}{\input{../cf_h2o/results/paper_artifacts/tables/target_information_ladder_table.tex}}
\captionof{table}{Target-information ladder. Zero-shot rows use no target transition labels before reporting; target offline adaptation rows use the stated fraction of passive target route labels and evaluate on held-out target routes except for the oracle row.}
\label{tab:target-information-ladder}
\end{center}

\subsection{Guarded full-module source ensemble}

The full-module diagnostic tests whether the rigid \ours residual can be improved without target-label calibration.
The best selected variant is a rigid-first source-gated delta with a guarded source-city ensemble.
It uses unlabeled target-city simulator summaries to weight source-city residual experts and to decide whether an ensemble is allowed; target transition labels remain held out.
The source-summary feature set is selected only by source-fold diagnostics, not by target performance; the selected feature sets are reported in \Cref{tab:guarded-source-feature}.
Across the six configured selector-stress splits, the guarded selector lowers mean total MSE from $884.853$ for the dense \Htwo-style baseline and $851.893$ for rigid \ours to $836.115$, with $6/6$ wins against \Htwo.
A same-information dense \Htwo source ensemble, using the same unlabeled target summaries for source weighting, improves to $846.466$ but remains weaker than the guarded mechanism-factored ensemble.
Because our dense baseline is a reproducible \Htwo-style residual-transfer surrogate rather than a full H2O+ reproduction, \Cref{tab:h2oplus-fidelity} lists the implemented and omitted components.

\begin{table}[t]
\centering
\scriptsize
\resizebox{\textwidth}{!}{\begin{tabular}{llllllllll}
\toprule
split & target\_city & selected\_candidate & selected\_is\_ensemble & max\_source\_weight & h2oplus\_mse & rigid\_cfcmt\_mse & guarded\_mse & guarded\_ratio\_vs\_h2oplus & guarded\_ratio\_vs\_rigid \\
\midrule
source\_open\_na\_to\_singapore & singapore lta & dense\_static\_local\_delta\_alpha\_0p75 & False & 0.465 & 2376.1427 & 2388.7628 & 2367.5165 & 0.9964 & 0.9911 \\
source\_singapore\_austin\_halifax\_to\_mbta & mbta & ensemble\_dense\_static\_local\_delta\_alpha\_0p50 & True & 0.5909 & 209.6939 & 116.4969 & 93.8321 & 0.4475 & 0.8054 \\
leave\_one\_city\_out\_all::singapore\_lta\_all & singapore lta & dense\_static\_local\_delta\_alpha\_0p75 & False & 0.465 & 2376.1427 & 2388.7628 & 2367.5165 & 0.9964 & 0.9911 \\
leave\_one\_city\_out\_all::austin\_capmetro\_all & austin capmetro & ensemble\_dense\_static\_local\_delta\_alpha\_0p25 & True & 0.5905 & 63.8731 & 47.6396 & 40.7936 & 0.6387 & 0.8563 \\
leave\_one\_city\_out\_all::halifax\_transit\_all & halifax transit & rigid\_cfcmt & False & 0.3749 & 73.5689 & 53.1982 & 53.1982 & 0.7231 & 1.0 \\
leave\_one\_city\_out\_all::mbta\_all & mbta & ensemble\_dense\_static\_local\_delta\_alpha\_0p50 & True & 0.5909 & 209.6939 & 116.4969 & 93.8321 & 0.4475 & 0.8054 \\
\bottomrule
\end{tabular}
}
\caption{Full-module guarded source-ensemble diagnostic. This table uses six configured selector-stress splits, including repeated target cities; it is therefore a source-selection diagnostic rather than the four-target strict leave-one-city-out count in \Cref{tab:strict-loo}.}
\label{tab:guarded-ensemble}
\end{table}

The guard is important.
MBTA and Austin have a dominant similar source under the selected target-summary descriptor, so the ensemble is allowed and improves the selected MSE substantially.
Halifax has diffuse source weights, so the guard falls back to rigid \ours and avoids the unguarded ensemble regression.
Per-mechanism hard splicing is weaker in this diagnostic, with mean total MSE $864.592$, because mechanism-level source-fold improvements do not reliably transfer to every target.
Additional same-information baselines, selector ablations, threshold sensitivity, target-label leakage checks, feature-set ablations, and oracle-gap diagnostics are reported in \Cref{tab:guarded-same-info,tab:guarded-selector-ablation,tab:guarded-threshold,tab:guarded-label-audit,tab:guarded-feature-ablation,tab:guarded-oracle-gap}.

\subsection{One-step policy proxy}

The strict one-step policy summary is more mixed but still favorable.
\ours improves reward/regret on three of four targets, improves mean absolute headway error on three of four, and lowers bunching rate on all four targets.
The mean regret ratio is $0.876$, and the mean bunching-rate ratio is $0.644$.
Halifax is the main policy exception: \ours lowers headway error slightly but has worse reward/regret than \Htwo.
This is consistent with a limited control-facing claim: mechanism-factored transfer improves cross-city dynamics strongly and often improves one-step policy-proxy quality, but it is not a closed-loop operational dominance result.

\begin{figure}[t]
\centering
\includegraphics[width=0.78\textwidth]{cross_city_policy_regret_ratio.png}
\caption{Strict leave-one-city-out one-step policy regret ratio. Ratios below 1 favor \ours over \Htwo.}
\label{fig:policy-regret}
\end{figure}

\begin{table}[t]
\centering
\scriptsize
\resizebox{\textwidth}{!}{\begin{tabular}{llllllllllllll}
\toprule
split & target\_city & h2oplus\_reward & cfcmt\_reward & cfcmt\_reward\_gain & cfcmt\_regret\_ratio & cfcmt\_headway\_error\_ratio & h2oplus\_bunching\_rate & cfcmt\_bunching\_rate & cfcmt\_bunching\_rate\_ratio & h2oplus\_large\_gap\_rate & cfcmt\_large\_gap\_rate & cfcmt\_large\_gap\_rate\_ratio & cfcmt\_hold\_seconds\_delta \\
\midrule
leave\_one\_city\_out\_all::singapore\_lta\_all & Singapore & -0.3338 & -0.3337 & 0.0001 & 0.8959 & 1.0007 & 0.0024 & 0.0024 & 0.9991 & 0.0016 & 0.0016 & 1.0 & -0.3597 \\
leave\_one\_city\_out\_all::austin\_capmetro\_all & Austin / CapMetro & -0.3326 & -0.332 & 0.0007 & 0.875 & 0.9995 & 0.0 & 0.0 & 0.0 & 0.0001 & 0.0 & 0.0 & 0.3909 \\
leave\_one\_city\_out\_all::halifax\_transit\_all & Halifax Transit & -0.3307 & -0.3312 & -0.0005 & 1.1113 & 0.9998 & 0.0 & 0.0 & 0.9355 & 0.0 & 0.0 & 1.0 & 1.4814 \\
leave\_one\_city\_out\_all::mbta\_all & MBTA & -0.3429 & -0.338 & 0.0049 & 0.6201 & 0.9967 & 0.0006 & 0.0004 & 0.6396 & 0.0008 & 0.0002 & 0.3174 & -1.1321 \\
\bottomrule
\end{tabular}
}
\caption{Strict leave-one-city-out one-step policy-proxy metrics. The table includes reward, regret ratio, headway error ratio, bunching, large-gap rate, and hold-time difference.}
\label{tab:policy}
\end{table}

\subsection{Source weighting is stable under a small grid}

The default source weighting uses $T=1.0$ and $\lambda=0.05$.
Across the $3 \times 3$ grid, all nine settings beat \Htwo on all six configured cross-city splits.
The mean weighted-\ours/\Htwo total-MSE ratio ranges only from $0.6025$ to $0.6149$.
The best setting in the grid is $T=0.5,\lambda=0$, but the default is not selected by post-hoc optimization and remains within the same narrow band.

\begin{figure}[t]
\centering
\includegraphics[width=0.58\textwidth]{source_weighting_sensitivity_heatmap.png}
\caption{Source-weighting sensitivity. Each cell reports mean total-MSE ratio versus \Htwo across configured cross-city splits.}
\label{fig:weight-sensitivity}
\end{figure}

\subsection{Target offline adaptation and mechanism-level errors}

The target-route offline-adaptation sweep tests whether a dense residual baseline with the same target passive-label budget quickly erases the no-target-label \ours advantage.
At $0\%$ target labels, weighted \ours has mean ratio $0.578$ against the dense source-only baseline.
When $1\%$, $5\%$, $10\%$, and $25\%$ of target routes are added to the dense baseline, the uncalibrated weighted-\ours/H2O ratio remains below $1$ on three of four strict targets, with mean ratios $0.581$, $0.586$, $0.617$, and $0.677$, respectively.
When \ours is allowed to use the same target route budget through the mechanism-weighted sufficient statistics, the corresponding mean ratios are $0.580$, $0.583$, $0.609$, and $0.660$.
The $100\%$ target-route point is an in-sample oracle reference and should not be interpreted as deployable no-calibration evidence.

\begin{figure}[t]
\centering
\includegraphics[width=0.70\textwidth]{calibration_budget_ratio.png}
\caption{Target-route offline-adaptation budget sensitivity. Ratios compare weighted \ours with no target labels or the same target route-label budget against a dense \Htwo-style baseline using that target budget.}
\label{fig:calibration}
\end{figure}

The same-information adaptation-baseline comparison gives a sharper design lesson.
AdaRL-style selected refit is nearly identical to full \ours refit because the target calibration slice usually selects all mechanism groups.
By contrast, the low-dimensional FANsRL-style context-bias adapter is better than full refit on average at every positive non-oracle budget, with mean ratios $0.559$, $0.545$, $0.572$, and $0.629$ against the dense adapted baseline.
It wins on four targets at the $1\%$ budget and on three of four targets at larger route budgets.
Thus the best few-shot variant is not unconstrained target refitting; it is mechanism-factored transfer plus a low-dimensional target context correction.

\begin{table}[t]
\centering
\scriptsize
\resizebox{\textwidth}{!}{\begin{tabular}{lllllllllll}
\toprule
budget & N & cal\_lines & cal\_trans & cfcmt\_refit & adarl\_refit & adarl\_bias & fansrl\_bias & fansrl\_wins & refit\_wins & oracle \\
\midrule
0.0 & 4 & 0.0 & 0.0 & 0.5782 & 0.5782 & 0.5782 & 0.5782 & 4 & 4 & False \\
0.01 & 4 & 3.5 & 3972.0 & 0.5801 & 0.5801 & 0.5585 & 0.5585 & 4 & 3 & False \\
0.05 & 4 & 23.0 & 30468.0 & 0.5831 & 0.5832 & 0.5472 & 0.5446 & 3 & 3 & False \\
0.1 & 4 & 50.5 & 61704.0 & 0.6094 & 0.6094 & 0.5718 & 0.5717 & 3 & 3 & False \\
0.25 & 4 & 132.75 & 167712.0 & 0.6596 & 0.6616 & 0.6318 & 0.6289 & 3 & 3 & False \\
1.0 & 4 & 534.75 & 697392.0 & 0.7242 & 0.7244 & 0.7222 & 0.7221 & 3 & 3 & True \\
\bottomrule
\end{tabular}
}
\caption{Same-information few-shot adaptation baselines. Ratios are total MSE relative to the dense \Htwo-style baseline using the same target route-label budget. The $1.0$ budget is an in-sample oracle reference.}
\label{tab:adaptation-baselines}
\end{table}

Per-mechanism errors show that the total-MSE advantage is not driven by one output group only.
Weighted \ours improves all three mechanism groups on average: headway/gap ratio $0.574$, demand/dwell/reward ratio $0.771$, and speed ratio $0.668$.
Headway/gap is the largest contributor to the main gain, which is consistent with the holding-control objective.

\begin{figure}[t]
\centering
\includegraphics[width=0.62\textwidth]{per_mechanism_error_ratio.png}
\caption{Per-mechanism mean MSE ratios under strict leave-one-city-out dynamics validation.}
\label{fig:mechanism-error}
\end{figure}

\subsection{Robustness and boundary conditions}

Generator robustness supports the main trend but also identifies failure modes.
Weighted \ours beats \Htwo on every city target in five of seven perturbation/noise scenarios.
Under headway-gap bias and mixed-mechanism bias, the result weakens to three of four targets.
Thus the evidence does not say that \ours is invariant to any target-generation change; it says that the result is robust to several plausible perturbations.

\begin{figure}[t]
\centering
\includegraphics[width=0.9\textwidth]{generator_robustness.png}
\caption{Generator-bias and measurement-noise robustness under strict leave-one-city-out evaluation.}
\label{fig:generator}
\end{figure}

Mechanism ablations support the mechanism-specific interpretation, with a caveat.
Wrong mechanism groupings and action/time-only features are much worse than \Htwo, with mean ratios above $7$.
Dense matched sparse features are close to full \ours, with mean ratio $0.772$, which means the result should be described as evidence for mechanism-wise sparse transfer rather than definitive causal identification.

The source-subset check is the strongest limitation.
With all four cities, weighted \ours has mean ratio $0.578$.
Excluding Austin remains favorable with mean ratio $0.663$.
Excluding Singapore makes the result nearly tied with \Htwo, with mean ratio $0.998$.
The two-city Halifax-MBTA pair is slightly worse than \Htwo, with ratio $1.042$.
This indicates that source diversity matters: the method should be framed as multi-source cross-city transfer, not as guaranteed transfer between every pair of cities.

\begin{figure}[t]
\centering
\includegraphics[width=0.9\textwidth]{source_subset_robustness.png}
\caption{Source-subset robustness. Removing data-source families reveals where the claim becomes weak.}
\label{fig:subset}
\end{figure}

Route-level bootstrap confidence intervals reinforce the same boundary.
Unweighted \ours has a route-bootstrap ratio confidence interval below $1$ for all four targets.
Weighted \ours has intervals below $1$ for Austin, Halifax, and MBTA, but Singapore crosses parity with interval $[0.993,1.002]$.
The conservative reading is that source weighting improves the mean strongly, while Singapore remains the least secure target.

\subsection{External real APC/AVL validation}

The public Austin raw APC validation is the strongest non-static check in the current paper.
Across two real APC/AVL datasets, the evaluation uses 11.57 million observed stop-event transitions after quality filtering.
Weighted \ours beats the dense \Htwo-style residual transfer baseline in both years: the total-MSE ratio is $0.330$ on July--December 2021 and $0.323$ on August 2016.
Unweighted \ours also beats the dense baseline, with ratios $0.528$ and $0.517$.
This supports the mechanism-factored transfer claim on real historical APC/AVL observations.

The result should still be read conservatively.
The APC-derived no-correction passive baseline has lower total MSE than both residual-transfer models on this real APC slice.
This does not mean that a free-running uncalibrated simulator is operationally better; it means that the baseline is given the observed current APC state and that source-only residual corrections can overshoot under real operational noise.
The overshoot is concentrated in headway/gap outputs, which dominate the unnormalized total MSE, while \ours is better than the passive baseline on several demand, dwell, and speed outputs.
Thus the real APC result strengthens the comparison between \ours and dense residual transfer, but it also shows that a safe residual gate or small target-city calibration budget is needed before claiming improvement over a state-informed passive baseline.

The few-shot APC sweep confirms this diagnosis.
With only $1\%$ of Austin line groups used for target residual calibration, the weighted-\ours residual gate improves over the passive baseline in both APC datasets, with ratios $0.970$ and $0.957$.
At $25\%$ target line-group budget, target-only \ours reaches the best ratios, $0.893$ on July--December 2021 and $0.879$ on August 2016.
Thus the real APC gap is not evidence against the mechanism-factored representation; it is evidence that source-only residual transfer needs target-domain trust calibration under real APC noise.
This is the same qualitative regime targeted by AdaRL/FANsRL-style few-shot adaptation, while the direct AdaRL/FANsRL codebases are environment-specific RL systems rather than drop-in bus residual evaluators.

The chronological time-budget sweep is stricter and closer to a deployment story.
One hour of target history is too small: only $2$ and $18$ calibration transitions are available in the two APC slices, and the residual gate remains above the passive baseline.
At six hours, the residual gate reaches parity or better in both datasets, with ratios $1.000$ and $0.961$.
At one day, source-plus-target refitting is below the passive baseline in both datasets, with ratios $0.943$ and $0.938$.
At one week, target-only mechanism-factored adaptation is best, with ratios $0.893$ and $0.876$.
This supports the intended target-offline-adaptation claim: a small but nontrivial amount of passive target history is enough to calibrate mechanism trust, whereas the pure source-only model should not be deployed without such a gate.

\subsection{Auxiliary sampled rollout}

The sampled BusSimEnv rollout is not primary evidence, but it checks whether the learned one-step action preferences are grossly inconsistent with executable line environments.
We expand this diagnostic to $80$ route episodes per policy across all four city bundles and add a threshold equalization baseline.
The result is mixed and weaker than the one-step policy proxy.
On raw mean reward, the dense \Htwo-style policy is best, with $-615.8$ versus $-644.9$ for weighted \ours.
However, raw rewards have different city scales.
After averaging city-level ratios, weighted \ours has lower headway error than \Htwo, with ratio $0.960$, and a similar bunching rate, with ratio $0.979$.
The city-stratified paired bootstrap gives a $95\%$ confidence interval of $[0.898,1.008]$ for the weighted-\ours/\Htwo headway-error ratio and $[0.877,1.052]$ for the bunching-rate ratio.
The city-level pattern remains uneven: \Htwo wins clearly on Austin reward, while weighted \ours improves headway or reward metrics on Halifax, MBTA, and Singapore.
Thus the expanded rollout should be read as a boundary condition, not as control superiority evidence: the dynamics-transfer result is strong, but closed-loop policy performance remains metric- and environment-sensitive.

We also add a libsumo full-line counterfactual diagnostic after inserting explicit bus-stop dwell points into generated SUMO scenarios.
The executable scenarios contain all validated lines in the four city bundles: $234$ Austin, $176$ Halifax, $940$ MBTA, and $789$ Singapore lines, for $58116$ generated bus-stop points.
The rollout uses a four-hour window, at most eight departures per line, the same $\{0,15,30,45,60\}$ second action grid for residual fitting and control, and excludes first-arrival warm-up events whose preceding headway is unobserved.
Across the reported policies, the diagnostic evaluates roughly $9.6$--$10.0$ thousand repeated-headway stop decisions after discarding about $7.0$--$8.0$ thousand warm-up stop events.
The unconstrained learned residual controllers remain weak in closed loop: dense \Htwo-MPC has mean headway error $685.55$, no holding has $683.84$, and weighted \ours-MPC improves to $645.35$ but still under-holds severe bunching cases.
Classical holding is much stronger on this generated topology: threshold equalization reaches $531.35$, and a retuned Daganzo rule with $\alpha=0.8$ reaches $529.64$.
The best current executable controller is a target-summary gated Daganzo prior.
It uses only target route/timetable summaries to choose a conservative or aggressive Daganzo gain and reaches mean headway error $528.63$, reward $-529.98$, and bunching rate $0.340$.
Against dense \Htwo-MPC, the city-stratified paired line bootstrap gives this target-summary prior a reward delta of $155.75$ with interval $[135.14,184.26]$, headway-error ratio $0.660$ with interval $[0.625,0.692]$, and bunching ratio $0.650$ with interval $[0.593,0.698]$.
This diagnostic changes the control-facing interpretation.
The mechanism-factored world model improves over dense learned residual MPC, but the closed-loop win over classical rules comes from contextualizing a physically meaningful control prior, not from free-form residual MPC.
The generated SUMO scenarios are still not field-calibrated counterfactuals, so the result is best read as a design constraint for the next controller: \ours should instantiate target mechanism state and gate a safe holding prior before allowing learned residual actions.

\subsection{RESCO traffic-signal stress test}

The traffic-signal stress test evaluates whether the same mechanism-transfer idea remains useful in a counterfactual control setting where strong classical rules are available.
We use eight RESCO SUMO scenarios, three matched random seeds, 600-second closed-loop rollouts, and the original SUMO \texttt{tlLogic} green phases as the action set.
We also include a generated native SUMO actuated-signal baseline that preserves the same network and route demand but lets SUMO's built-in actuated controller run without learned phase overrides.
The reported metric is mean queue averaged over target scenarios and seeds, with throughput ratio as a secondary check.

\begin{center}
\includegraphics[width=0.78\textwidth]{traffic_signal_resco_extended_summary.png}
\captionof{figure}{Extended RESCO multi-seed phase-transfer validation. Bars show mean queue across eight held-out RESCO target scenarios and three matched seeds; error bars show seed-level standard deviation of mean queue.}
\label{fig:resco-traffic-signal}
\end{center}

\begin{center}
\scriptsize
\resizebox{\textwidth}{!}{\begin{tabular}{lrrrrrrr}
\toprule
Policy & Mean queue & Seed std & P90 queue & Trip time & Trip wait & Time loss & Throughput \\
\midrule
MaxPressure & 12.9068 & 0.1546 & 19.5279 & 87.5891 & 8.0605 & 24.5171 & 0.8267 \\
CFCMT + pressure guard & 12.9180 & 0.1919 & 19.5973 & 87.4863 & 8.0114 & 24.4127 & 0.8256 \\
Phase pressure & 13.2446 & 0.1816 & 19.3911 & 86.5371 & 7.4544 & 23.4628 & 0.8199 \\
H2O+-style dense & 17.2279 & 3.6084 & 28.9325 & 86.6585 & 8.5768 & 24.5488 & 0.7653 \\
Target-static simulator & 20.4023 & 0.1758 & 32.5538 & 97.9633 & 19.9032 & 35.3928 & 0.7429 \\
Native actuated SUMO & 30.1825 & 0.3536 & 47.6873 & 119.7813 & 33.0017 & 57.0805 & 0.7580 \\
Fixed RESCO program & 42.6732 & 0.2709 & 69.6852 & 142.3116 & 54.6333 & 79.8759 & 0.7127 \\
\bottomrule
\end{tabular}
}
\captionof{table}{Extended RESCO traffic-signal stress test. Trip time, trip wait, and time loss are official SUMO \texttt{tripinfo} metrics over completed vehicles and should be read together with throughput; lower queue and higher throughput are primary.}
\label{tab:resco-traffic-signal}
\end{center}

\Cref{fig:resco-traffic-signal,tab:resco-traffic-signal} show the main result.
MaxPressure has mean queue $12.907$, while the pressure-guarded CFCMT phase controller has mean queue $12.918$.
Their official \texttt{tripinfo} metrics are essentially tied: CFCMT has trip time $87.486$, trip wait $8.011$, and time loss $24.413$, compared with $87.589$, $8.061$, and $24.517$ for MaxPressure.
The next strongest classical row is phase pressure with mean queue $13.245$; it has lower trip time, wait, and time loss, but also lower throughput, so these trip metrics must be interpreted with queue and throughput rather than in isolation.
The dense \Htwo-style phase residual, target-static simulator MPC, generated native actuated SUMO baseline, and fixed RESCO programs are weaker, with mean queues $17.228$, $20.402$, $30.183$, and $42.673$, respectively.
The paired seed-target bootstrap for guarded CFCMT versus MaxPressure has mean delta $0.011$ with 95\% interval $[-0.046,0.079]$, so the correct conclusion is near-parity with MaxPressure rather than a significant win.
Against phase pressure the mean delta is $-0.327$ but the interval $[-0.759,0.023]$ still crosses zero.
Against dense \Htwo-style residual transfer, target-static simulator MPC, native actuated SUMO, and fixed programs, the intervals favor guarded CFCMT: the mean queue deltas are $-4.310$, $-7.484$, $-17.265$, and $-29.755$, respectively.
Per-target winners remain mixed, which is important for claim scope: several regular or sparse networks are best handled by MaxPressure or a simulator/pressure variant, while the guarded CFCMT controller avoids the large negative transfer seen in dense residual and simulator-only MPC.

The 3600-second long-horizon check in \Cref{tab:resco-traffic-signal-3600s} is stricter and exposes a boundary.
Phase pressure is best at mean queue $32.234$ and also has the best trip time, $104.229$.
MaxPressure and guarded CFCMT rise to mean queues $42.787$ and $43.511$, with trip times $106.918$ and $108.633$.
Guarded CFCMT still beats dense \Htwo-style residual transfer, native actuated SUMO, and fixed programs, with mean queue deltas $-5.995$, $-14.836$, and $-41.098$ under the paired seed-target bootstrap, but it is significantly worse than phase pressure on this horizon.
Thus the stress test supports guarded transfer as a safe residual design, not as a long-horizon replacement for a well-tuned pressure rule.

This result changes the control-facing message.
The evidence is not that CFCMT universally dominates classical traffic control.
It is that a mechanism-transfer model can be coupled with a source-only pressure guard to reach MaxPressure-level closed-loop safety on a broad RESCO subset while substantially improving over dense residual transfer and simulator-only MPC under the same phase-transfer protocol.

\paragraph{Native-protocol RL cross-checks.}
\Cref{tab:official-resco-rl-full,tab:libsignal-rl-full} report full-budget external RL runner results.
The official RESCO adapter completed 42 supported runs with zero failures: \texttt{IDQN} on all eight extended scenarios and upstream-supported \texttt{MPLight} on six scenarios, while the known upstream \texttt{MPLight} crashes on \texttt{cologne8} and \texttt{ingolstadt21} are excluded rather than patched into a different algorithm.
The patched LibSignal adapter completed 36/36 SUMO runs with zero failures across \texttt{DQN}, \texttt{PressLight}, \texttt{FRAP}, and \texttt{MPLight} on \texttt{sumo1x1}, \texttt{sumo1x3}, and \texttt{sumo4x4}.
These tables strengthen the experimental audit trail, but they should not be read as direct wins or losses against \ours: RESCO and LibSignal use their own observations, rewards, action timing, and metrics.
The same-protocol conclusion remains the one above: guarded \ours improves over dense residual transfer and simulator-only MPC and reaches pressure-rule-level safety on the short RESCO stress test, while phase pressure is the stronger long-horizon classical controller.
